\documentclass[aps,prd,twocolumn,superscriptaddress,longbibliography,nofootinbib,floatfix,10pt]{revtex4-2}

\usepackage{amsmath,amssymb}
\usepackage{graphicx}
\usepackage{booktabs}
\makeatletter
\let\ACM@PackageWarning\PackageWarning
\def\PackageWarning#1#2{%
  \def\ACM@warnpkg{#1}%
  \def\ACM@namerefpkg{nameref}%
  \ifx\ACM@warnpkg\ACM@namerefpkg
  \else
    \ACM@PackageWarning{#1}{#2}%
  \fi
}
\usepackage{hyperref}
\let\PackageWarning\ACM@PackageWarning
\makeatother
\usepackage{iftex}
\usepackage{xcolor}
\usepackage{tikz}
\usetikzlibrary{arrows.meta,positioning}
\ifXeTeX
  \usepackage{fontspec}
  \setmainfont{Times New Roman}
  \IfFontExistsTF{Microsoft YaHei}{
    \newfontfamily\zhnamefont{Microsoft YaHei}
  }{
    \newfontfamily\zhnamefont{SimSun}
  }
  \newcommand{\zhname}{{\zhnamefont 刘海涛}}
\else
  \newcommand{\zhname}{刘海涛}
\fi

\begin{document}

\title{ACM Forensic Audit of MOND-Favored Galaxies:\\Residual Pathologies in the SPARC Sample}

\author{Liu Haitao / \zhname}
\affiliation{Independent Researcher, Shanghai, China}
\email{lesleyliuhaitao@hotmail.com}

\date{\today}

\begin{abstract}
We present an ACM-based forensic audit of the galaxies that appear to favor
MOND in the SPARC sample. The purpose of the paper is not to introduce yet
another replacement function or empirical low-acceleration curve. It is to ask
whether the MOND-favored regime represents a unified, stable, and reproducible
physical mechanism, or whether it is better understood as a mixed residual
regime shaped jointly by observational and modeling conventions. Paper II
therefore adopts the retained ACM galaxy trunk established in Paper I and
shifts the focus from global model comparison to full-sample residual
pathology in SPARC~\cite{LelliMcGaughSchombert2016SPARC}.

Under a single ACM baseline applied to 164 comparison objects, the ACM-better
systems form the dominant physical cloud, whereas the MOND-favored minority
does not behave like a coherent dynamical counterexample. Instead, it separates
into three distinct mechanisms: a geometry-sensitive subset that flips under
distance-edge stress tests; a stellar-normalization-sensitive subset recoverable
under conservative stellar mass-to-light rescaling; and a hard31 subset whose
outer gas distributions are unusually flat, extended, and structure-poor, i.e.
structure-poor outer-disk cases. This
decomposition suggests that the MOND-favored regime is better interpreted as a
set of residual mechanisms stabilized under historical limitations, rather than
as a unified physical stronghold.

The claim of this paper is not that earlier researchers committed deliberate
errors. MOND's historical success was real and important as an empirical
low-acceleration framework~\cite{Milgrom1983MOND,FamaeyMcGaugh2012}. Our claim is
instead that, under earlier distance systems, stellar-normalization practices,
and low-structure outer-disk regimes, a local threshold law could absorb
multiple residual structures and thereby stabilize an interpretively
overfit-friendly observational regime. We do not declare a final verdict on
behalf of future data. We argue, however, that harder distance anchors, more stable stellar
mass normalization, and higher-resolution outer-disk measurements will
continue to arbitrate the interpretation of these regimes. The MOND-favored
regime is thus revealed as a collection of residual mechanisms stabilized under
historical observational conventions, rather than a coherent physical domain.
\end{abstract}

\maketitle

\begin{center}
\fcolorbox{black}{gray!10}{%
\parbox{0.92\linewidth}{%
\textbf{Interactive Audit Interface.}
An interactive visual companion to Paper~II is available at
\href{https://lesleyliuhaitao-ctrl.github.io/galaxy-audit-system/}{lesleyliuhaitao-ctrl.github.io/galaxy-audit-system/}.}}
\end{center}

\section{Introduction}

\subsection{Starting from the unified trunk of Paper I}

The ACM galaxy-dynamics trunk is not introduced here for the first time.
Paper I defined the retained galaxy trunk, including the unified constant
family, the local background-response chain, the propagation-layer logic, and
the global matched-sample performance of the stable trunk~\cite{Liu2026PaperI}\footnote{Paper I preprint: \href{https://doi.org/10.5281/zenodo.19163963}{Zenodo DOI}. Code repository: \href{https://github.com/lesleyliuhaitao-ctrl/anchor-collapse-model.git}{anchor-collapse-model}.}. More
importantly, Paper I fixed the methodological discipline that Paper II
continues to follow: failed objects are retained rather than erased by
per-galaxy patches, and explanatory power is not purchased by post hoc degrees
of freedom.

Paper II therefore does not take on the task of re-deriving the first
principles of the ACM galaxy trunk. All ACM calculations, residual
classifications, and regime audits in this paper are performed on top of the
single baseline already established in Paper I. The question here is not
whether ACM works globally in the first place. The follow-up question is more
specific: once a unified ACM baseline has been fixed, what does it mean that a
minority of systems still appears to favor MOND? Do these systems form a
coherent physical counterexample, or do they instead compress several distinct
bias channels and fragile regimes into what only looks like a stable MOND
stronghold?

\subsection{From model comparison to residual pathology}

Comparisons between ACM and MOND can easily remain at a coarse level: which
model has the lower median residual, which model wins on more galaxies, and
which model follows the low-acceleration empirical curve more closely. Such
comparisons are useful, but they do not answer the specific question pursued in
this paper. Once a minority of systems appears MOND-favored, repeating that ACM
is globally stronger does not by itself explain what these galaxies actually
represent.

This paper therefore shifts the object of study from simple model competition
to residual pathology. We do not treat the MOND-favored subset as a naturally
unified counterexample class. We treat it instead as a set of residual regimes
that must be audited: are these systems geometry-locked, stellar
normalization-biased, or so flat in their outer gas structure that any trunk
relying on local structural activation is driven into an information-poor
regime? In that sense, the central product of this paper is a full-sample
Residual Pathology Map rather than a ranked win--loss table. This shift is motivated
not only by ACM's retained baseline, but also by the observational framing
established by SPARC~\cite{LelliMcGaughSchombert2016SPARC}, MOND
phenomenology~\cite{FamaeyMcGaugh2012}, and the radial-acceleration
literature~\cite{McGaughLelliSchombert2016RAR}. In practical terms, Paper II
uses the same high-quality SPARC sample provided by
Ref.~\onlinecite{LelliMcGaughSchombert2016SPARC}, while stripping away the
interpretive layers associated with the MOND phenomenology and
radial-acceleration reading of that sample.

This shift carries a deeper consequence. The aim here is not to put MOND on
trial as a historical actor. MOND's historical success was real. It provided a
courageous local approximation and, for a long period, served as a bridge
between low-acceleration observations and theoretical tension. What this paper
audits is how that local success became stabilized under historical
conditions, and how its structural limits emerge once the sample is examined at
higher residual resolution. It is in this restricted methodological sense that
we use the word \emph{overfitting}: not to suggest subjective manipulation, but
to describe a regime-dependent structure stabilized under historical
observational conditions.

\section{Audit Philosophy and Unified ACM Baseline}

\subsection{Why a unified baseline matters}

All analyses in this paper are built on the stable galaxy trunk established in
Paper I~\cite{Liu2026PaperI}. In practice this means that we do not refit ACM
freely on a galaxy-by-galaxy basis, nor do we introduce new free parameters in
order to rescue a small stubborn subset. We deliberately place ACM in a strict,
unified, and reproducible audit posture: one trunk function, one constant
family, and one local-response chain applied across the full comparison sample.

The purpose of this choice is not simply parameter economy. It is to prevent
Paper II from collapsing back into empirical curve chasing. The goal is not to
show that ACM can also behave like an old-style fitter with enough targeted
adjustments. Rather, the goal is to test how stable the purported MOND-favored
regime really is. If a subset remains MOND-favored only so long as the distance
scale stays soft, the stellar mass normalization stays flexible, or the source
family remains within a particular observational regime, then that apparent
loyalty to MOND must itself become an audit object.

\subsection{The four audit actions}

To pursue this question, Paper II uses four mutually reinforcing audit actions.
The first is the \emph{distance-edge test}, in which the adopted distance is
shifted within its quoted uncertainty, $D \pm e_D$, and the verdict response is
measured. The second is the \emph{stellar normalization scan}, in which the
stellar mass-to-light ratio is varied within a conservative and astrophysically
interpretable range. The third is the \emph{outer-gas-shape diagnostic}, which
uses outer-disk gas slope, curvature, and persistence to identify systems whose
local structural signal is intrinsically weak. The fourth is the \emph{source-
and-family audit}, which tests whether distance-source families, reference
chains, and observational regimes couple systematically to the residual
mechanisms. These audit actions are motivated by the status of SPARC mass
models~\cite{LelliMcGaughSchombert2016SPARC}, the LSB rotation-curve
literature~\cite{McGaughRubinDeBlok2001}, stellar mass-normalization
practice~\cite{BellDeJong2001}, and the distinction between soft and hard
distance anchors such as the TRGB method~\cite{MadoreMagerFreedman2009}.

All four actions share the same principle: none of them attempts to smuggle a
new free parameter into the ACM trunk. The audit target is the sample and the
interpretive regime, not the invention of a new decorative patch.

The comparison sample itself is not newly defined in Paper II. The 164-object
set used throughout this audit is the same matched comparison set inherited
from Paper I, with the same inclusion criteria and the same object-level
baseline used for direct ACM--MOND comparison. Paper II therefore changes the
mode of analysis, not the underlying sample definition. The released
object-level table for this inherited comparison set is the Paper II audit
export \texttt{full\_sample\_residual\_pathology\_audit.csv}.

\subsection{Operational distance-quality tiers}

The geometry audit uses the SPARC distance flag $f_D$ as an operational proxy
for distance quality, but not as a replacement for the original SPARC metadata.
For the purposes of Paper II, we group the flag values into three audit tiers.
The \emph{soft or derived tier} consists of $f_D \leq 2$, where the adopted
distance remains comparatively exposed to derived or soft priors. The
\emph{intermediate tier} consists of $f_D = 3$, which is rare in the present
comparison sample. The \emph{hard-anchor tier} consists of $f_D \geq 4$, with
$f_D = 5$ treated as the strongest geometric-anchor class in the audit. This
is an operational audit mapping designed to test regime stability under
distance quality; it is not a redefinition of the SPARC catalog itself. In
practical observational terms, the hard-anchor language is intended to point
toward TRGB-, Cepheid-, maser-, or comparably rigid geometric-distance
regimes, whereas the soft/derived tier captures the classes of distance inputs
most exposed to indirect or historically softer calibration chains.

\subsection{From 62 to 22 / 9 / 31}

The decomposition of the MOND-favored minority was not obtained by subjective
inspection. It emerged from a clear audit sequence. First, the unified trunk
was applied to the 164-object comparison sample, yielding an initial split into
102 ACM-better systems and 62 MOND-favored systems. Second, the 62-object
subset was subjected to the distance-edge test. Twenty-two of these systems
flip back to ACM under edge shifts, and the flips occur overwhelmingly on the
$D + e_D$ side. The first audit layer therefore removes a geometry-sensitive
subset from the supposed MOND stronghold. The full object-level audit table for
this inherited comparison set is released with the evidence package associated
with Paper II.

Third, the remaining holdout systems were subjected to a conservative stellar
mass-normalization scan. Within the residual 40-object pool, nine systems
return to ACM under plausible $M/L$ adjustments. Only after these first two
filters does the final hard31 subset emerge. The three mechanisms discussed in
this paper are therefore not preconceived labels; they are the outcome of a
stepwise sequence of distance-prior audit, stellar-normalization audit, and
outer-structure/source-regime audit. The procedural chain is summarized in
Fig.~\ref{fig:audit-flow}, while the resulting full-sample decomposition is
visualized in the Residual Pathology Map of Fig.~\ref{fig:pathology-map}.

\begin{figure*}[t]
\centering
\begin{tikzpicture}[
  node distance=1.15cm and 0.9cm,
  every node/.style={font=\small, align=center},
  box/.style={draw=black!70, rounded corners=4pt, line width=0.7pt, fill=black!2, inner sep=8pt, text width=0.21\textwidth},
  verdict/.style={draw=blue!50!black, rounded corners=4pt, line width=0.8pt, fill=blue!4, inner sep=8pt, text width=0.19\textwidth},
  arrow/.style={-{Latex[length=3mm,width=2mm]}, line width=0.8pt, draw=black!70}
]
\node[box] (sample) {164 comparison objects\\under the retained ACM baseline};
\node[box, right=of sample] (distance) {Distance-edge audit\\$D \pm e_D$\\distance-quality screening};
\node[verdict, right=of distance] (geom) {22 geometry hostages\\removed from the MOND-favored pool};
\node[box, below=of distance] (ml) {Conservative stellar-normalization audit\\bounded $M/L$ rescaling};
\node[verdict, right=of ml] (stellar) {9 stellar hostages\\returned to the ACM baseline};
\node[verdict, left=of ml] (hard31) {31 hard31\\structure-poor outer-disk cases};

\draw[arrow] (sample) -- (distance);
\draw[arrow] (distance) -- (geom);
\draw[arrow] (distance) -- (ml);
\draw[arrow] (ml) -- (stellar);
\draw[arrow] (ml) -- (hard31);
\end{tikzpicture}
\caption{
Sequential audit flow used in Paper II. The decomposition of the MOND-favored
minority emerges through explicit filtering stages rather than through post hoc
relabeling: distance-edge audit first isolates the geometry-hostage subset,
conservative stellar-normalization audit then isolates the stellar-hostage
subset, and the residual hard31 regime remains as the final structure-poor
audit target.
}
\label{fig:audit-flow}
\end{figure*}

\subsection{Minimal mathematical definitions}

Although this paper does not re-derive the ACM trunk, the audit procedure must
remain mathematically reproducible. All comparisons in this paper inherit the
unified retained trunk of Paper I, whose dynamical closure is written as
\begin{equation}
g_{\mathrm{ACM}}(r)=g_{\mathrm{bar}}(r)+\frac{2}{3}c^2\eta_{\mathrm{local}}(r),
\end{equation}
where $g_{\mathrm{bar}}(r)$ is the baryonic acceleration and
$\eta_{\mathrm{local}}(r)$ is the retained local background-response field from
Paper I~\cite{Liu2026PaperI}. Paper II does not repeat the first-principles
derivation of $\eta_{\mathrm{local}}$; it inherits it as part of the single
audit baseline.

On top of that baseline, only a few additional mathematical definitions are
introduced. The distance-edge test is implemented as
\begin{equation}
D_{\mathrm{new}} = D \pm e_D,
\end{equation}
with radius and baryonic components rescaled consistently under the shifted
distance. The stellar normalization audit is written as
\begin{equation}
\Upsilon_\star \rightarrow s\,\Upsilon_\star,
\end{equation}
where $s$ is restricted to a conservative and astrophysically interpretable
range. In the present audit this window is the same bounded envelope used in
the derived evidence tables, namely $s \in \{0.7, 0.85, 1.0, 1.15, 1.3\}$.

For the hard31 structural audit, the relevant quantities are not new dynamical
equations but diagnostic structure measures. The outer-disk gas slope is
defined by
\begin{equation}
S_{\mathrm{gas}}(r)=\left|\frac{d\ln(V_{\mathrm{gas}}^2)}{d\ln r}\right|,
\end{equation}
and the outer-disk gas curvature by
\begin{equation}
C_{\mathrm{gas}}(r)=\left|\frac{d^2\ln(V_{\mathrm{gas}}^2)}{d(\ln r)^2}\right|.
\end{equation}
These quantities are not treated as new corrective theory terms. They are used
as structural fingerprints that distinguish systems with genuinely weak local
structure from systems that remain information-rich but dynamically difficult.

\section{Global Residual Pathology Map}

\subsection{The 164-object pathology atlas}

To move from anecdotal cases to sample-level structure, the full 164-object
comparison sample is projected into a residual pathology map. The result is
clear. The ACM-better 102 do not form a marginal survival group; they form the
dominant physical cloud of the sample. By contrast, the MOND-favored minority
does not behave like a single physical family. It separates into three
distinct mechanisms.

The geometry-hostage 22 occupy the high-distance-error corner and are enriched
in soft distance sources. The stellar-hostage 9 inhabit the low-stellar-light,
gas-rich sector, but can be restored by conservative stellar mass
normalization. The gas-flat hard31 form a low-structure regime characterized by
flat outer gas profiles, weak curvature, high outer persistence, and the
absence of hard geometric distance anchors. In simple numerical terms, the
sample becomes 164 objects decomposed into 102 ACM-better systems and 62
apparent MOND-favored systems, which in turn separate into 22 geometry
hostages, 9 stellar hostages, and 31 hard31 judges---the latter being
structure-poor outer-disk cases that survive the first two audit layers. A
subsequent rigid-trunk audit leaves this 102/22/9/31 layering unchanged,
showing that the pathology map is not an artifact of a softer engineering
realization.

\begin{figure*}[!t]
\centering
\includegraphics[width=0.95\textwidth]{../../research_assets/derived_exports/full_sample_residual_pathology_map.png}
\caption{
Residual Pathology Map for the 164-object comparison set. The
ACM-better cloud serves as the primary physical trunk, while the three
MOND-favored residual regimes---geometry hostages, stellar hostages, and
gas-flat hard31---are isolated for forensic auditing.
}
\label{fig:pathology-map}
\end{figure*}

\subsection{A four-quadrant verdict}

The four-quadrant pathology map yields a clear verdict. In the geometry plane,
the geometry-hostage 22 cluster in the high-error sector and are the only class
that shows a systematic verdict reversal under $D+e_D$. In the mass/baryon
plane, both the stellar-hostage 9 and the gas-flat hard31 occupy
low-stellar-light, gas-rich territory, but only the former return under
conservative $M/L$ recovery. In the outer-gas-shape plane, the hard31
concentrate in the low-slope, low-curvature, high-persistence corner. They are
not gas-free systems. They are systems whose gas support is too flat, too
extended, and too poor in local structure to provide strong local feedback.

The importance of the pathology map is that it dissolves the MOND-favored
region into multiple residual mechanisms. The subsequent sections are therefore
not arbitrary case studies; they are successive unfoldings of the full-sample
pathology map.

\section{Geometry Hostages}

The distance-edge test produces one of the sharpest signals in the paper.
Among the 62 MOND-favored systems, 22 flip to ACM once the distance is moved to
the edge of the quoted uncertainty range. These reversals occur overwhelmingly
on the $D + e_D$ side rather than symmetrically on both sides. The median
relative distance error of this subset is approximately $24.96\%$, and its
internal fraction of $f_D=1$ systems is about $77.3\%$, both substantially
above the level of the dominant sample cloud.

\begin{table}[t]
\centering
\small
\caption{
Operational distance-quality composition of the geometry-hostage subset
compared with the ACM-better majority. Here ``soft/derived'' denotes
$f_D \leq 2$, ``intermediate'' denotes $f_D = 3$, and ``hard-anchor''
denotes $f_D \geq 4$.
}
\label{tab:distance-quality}
\begin{tabular}{lrrrr}
\toprule
Subset & $N$ & Soft/derived & Intermediate & Hard-anchor \\
\midrule
Geometry hostages & 22 & 20 & 0 & 2 \\
ACM-better cloud & 102 & 73 & 2 & 27 \\
hard31 & 31 & 30 & 0 & 1 \\
\bottomrule
\end{tabular}
\end{table}

These signals support a simple interpretation: the geometry-hostage 22 are not
natural MOND victories. They are systems whose verdict is being locked by the
current geometric scale. Their apparent MOND advantage does not require a new
dynamical regime. It requires only a slightly compressed distance prior capable
of holding a local threshold law in an artificially favorable position. In that
sense, they are better interpreted as geometry problems than as dynamics
problems. In operational distance-quality terms, the geometry-hostage subset is
concentrated overwhelmingly in the soft or derived tier, contains no
$f_D = 3$ intermediates, and contains no $f_D = 5$ strongest-anchor systems at
all. The object-level distance-quality evidence remains traceable in the
inherited comparison metadata released with the Paper II evidence package. The
geometry-hostage cluster itself is isolated in
Fig.~\ref{fig:recoverable-subsets}.

\section{Stellar Hostages}

The second audit layer reveals a smaller subset with strong sensitivity to
stellar mass normalization. When the remaining holdout systems are scanned
under conservative $M/L$ rescaling, 9 of the 40 return to ACM. The most common
recovery scale appears on the higher but still plausible end of the allowed
range. These systems therefore do not require a new dynamical law. They behave
instead like systems that have been held on the MOND-favored side by a weak
stellar normalization anchor, consistent with the long-standing ambiguity in
disk-galaxy mass-to-light estimates~\cite{BellDeJong2001,LelliMcGaughSchombert2016SPARC}.

\begin{table}[t]
\centering
\small
\caption{
Recovered stellar-hostage systems under conservative stellar
mass-to-light rescaling. All nine recovered objects return to the ACM baseline
at the same upper conservative scale, $s = 1.3$.
}
\label{tab:stellar-recovery}
\begin{tabular}{lrrc}
\toprule
Galaxy & Baseline $\Delta\mathrm{CPP}$ & Best $s$ & Returns to ACM \\
\midrule
NGC3741  & -7.82   & 1.3 & Yes \\
F574-1   & -9.82   & 1.3 & Yes \\
UGC12506 & -10.13  & 1.3 & Yes \\
NGC1705  & -25.51  & 1.3 & Yes \\
UGC00731 & -27.33  & 1.3 & Yes \\
NGC0247  & -43.22  & 1.3 & Yes \\
NGC2841  & -69.18  & 1.3 & Yes \\
NGC7814  & -81.39  & 1.3 & Yes \\
DDO154   & -401.51 & 1.3 & Yes \\
\bottomrule
\end{tabular}
\end{table}

This class shares a low-stellar-light, gas-rich background with the hard31, but
it does not share the same pathology. The stellar-hostage 9 are not entirely
structure-poor. Their key feature is that they are especially sensitive to the
normalization of the stellar component. They are therefore closer to
recoverable normalization bias than to a genuinely separate low-acceleration
dynamical regime. Importantly, the stellar-normalization audit is not a free
parameter search. It is a bounded normalization stress test: the recovery is
restricted to a conservative stellar-population envelope and, in the present
sample, the recovered cases all converge to the same upper conservative factor
rather than scattering across arbitrary values. Figure~\ref{fig:recoverable-subsets}
shows the recovery trajectories, and Table~\ref{tab:stellar-recovery} records
the object-level verdict return.

\begin{figure*}[t]
\centering
\begin{minipage}[t]{0.48\textwidth}
\centering
\includegraphics[width=\textwidth]{../../research_assets/derived_exports/pathology_geometry_hostages_focus.png}
\vspace{0.4em}

\small\textbf{Left: Geometry-hostage audit.}
Focused audit of the 22 systems that return to ACM under distance-edge stress
tests. The systematic flip on the $D+e_D$ side identifies these objects as
geometric artifacts rather than dynamical outliers.
\end{minipage}
\hfill
\begin{minipage}[t]{0.48\textwidth}
\centering
\includegraphics[width=\textwidth]{../../research_assets/derived_exports/pathology_stellar_hostages_focus.png}
\vspace{0.4em}

\small\textbf{Right: Stellar-hostage audit.}
Recovery trajectory for the 9 stellar-normalization-sensitive systems. Orange
points mark the initial MOND-favored states; blue triangles denote their
return to the ACM baseline after conservative $M/L$ auditing.
\end{minipage}
\caption{
Sequential audit of the recoverable MOND-favored subsets. The left panel
isolates the geometry-sensitive systems removed by the distance-edge test,
whereas the right panel isolates the stellar-normalization-sensitive systems
recovered under conservative $M/L$ auditing.
}
\label{fig:recoverable-subsets}
\end{figure*}

\section{The hard31 as Internal Audit Judges}

The status of the gas-flat hard31 is categorically different from that of the
first two groups. These 31 systems were not initially selected as ammunition
against MOND. They were first used as an internal ACM audit benchmark. In
other words, the first suspicion was not that the data were biased, but that
the ACM trunk might still be missing a sufficiently natural, as-yet-unmerged
correction.

For precisely that reason, the hard31 became the most severe test set. A series
of physically natural and zero-new-parameter operator variations was explored:
first-gradient schemes, second-curvature schemes, diffusion accumulation terms,
self-shielding terms, and nonlinear impedance couplings. The associated raw
experimental notes and archived negative trials will be released
progressively as research records at an appropriate stage. The issue is not
that these ideas were directionless. The issue is that none of them
materially changed the fate of the hard31. They were either too weakly
activated or only produced small perturbations. None was sufficient to bring
the systems back across the verdict boundary.

More importantly, the hard31 exhibit a coherent set of quantitative
fingerprints. Their number of hard geometric anchors is zero ($f_D=5=0$). Their
mean outer gas slope is approximately $1.12$, their mean outer gas curvature is
approximately $3.96$, and their mean outer-to-inner gas persistence ratio is
approximately $5.95$. In mechanical terms, such systems are close to an
``idle-running'' state: the outer disk is so flat, extended, and poor in
structure that insufficient structural variation remains to generate stable
feedback. Any model that depends on local-response activation will therefore
enter a weak-coupling regime here. For that reason, the hard31 do not resemble
a simple pile of bad points. They instead define a class of soft-source,
low-anchor, low-structure systems occupying a special observational regime. In
that limited but important sense, the hard31 behave like a featureless regime:
any sufficiently smooth low-acceleration curve can appear successful there,
but none is structurally necessary.

This is the turning point of the paper. Once all naturally favorable ACM-side
operator modifications fail systematically, the problem is forced to move from
``what function is still missing from the model'' to ``whether the input system
itself carries historical bias.'' This is why the hard31 deserve the name
\emph{judges}. They are neither MOND's hostages nor ACM's enemies. They are the
methodological subset that forces the audit to move from trunk-side ornament to
input-side diagnosis.

The positive role of the hard31 is reinforced by their source topology. Their
reference-family audit shows that they do not inhabit a hard geometric control
group. They remain concentrated in softer source regimes and lack hard distance
anchors entirely. This does not by itself prove input-side bias, but it does
establish a strong regime-level audit concern. In that sense, hard31 is better
treated not as a clean MOND fortress but as a soft-source observational island
that must remain in the main text rather than being relegated to supplementary
material. The relevant structural and source-topology evidence is shown in
Figs.~\ref{fig:hard31-structure} and~\ref{fig:hard31-topology}.

\begin{figure*}[t]
\centering
\includegraphics[width=0.92\textwidth]{../../research_assets/derived_exports/pathology_hard31_focus.png}
\caption{
Pathological fingerprinting of the hard31 subset. The subpanels isolate the
low-slope, high-persistence outer-gas regime, demonstrating that these systems
occupy a structure-poor observational blind spot.
}
\label{fig:hard31-structure}
\end{figure*}

\begin{figure*}[!t]
\centering
\includegraphics[width=\textwidth]{../../research_assets/derived_exports/hard31_reference_topology.png}
\caption{
Distribution of the hard31 subset within source-regime space. The concentration
in softer reference families and the lack of hard geometric anchors support the
interpretation of hard31 as a data-regime island rather than a unified physical
counterexample.
}
\label{fig:hard31-topology}
\end{figure*}

\section{Discussion}

Taken together, the results support a restrained but sharp reinterpretation:
the MOND-favored regime should no longer be discussed as a unified physical
stronghold. It behaves more like a distance--mass--model trap, or more broadly
like an overfit-friendly observational regime stabilized under historical
conditions.

This interpretation should not be read as an accusation of deliberate error.
MOND's historical success was real. Under the observational conditions of its
time, it provided a highly insightful local approximation and helped drive the
field's re-engagement with low-acceleration phenomena. The point of the present
audit is different. When the distance source is softer, the stellar mass
normalization more flexible, and the outer disk flatter and poorer in local
structure, a local threshold law can absorb several distinct residual
structures and compress them into an apparently unified success region.

In that sense, \emph{overfitting} is not used here as a moral category. Nor is
it meant in the narrow sense of a merely over-parameterized fit. It is a
methodological diagnosis at the interpretive and regime level. It describes a
structure in which multiple residual regimes that should not belong to the same
dynamical category become compressed, aligned, and stabilized under historical
observational conditions. The audit of the 22, the 9, and the 31 is precisely
the process of decomposing this stabilized overfit regime.

For hard31-like systems, the next discriminant is therefore not another
decorative local function, but a harder observational test. If future data add
stronger geometric anchors, higher-resolution outer-disk structure, and more
stable inclination fields, then these systems will no longer be able to remain
ambiguous on the basis of low-structure smoothness alone. The relevant
question is not whether a smoother empirical curve can still be drawn, but
whether the regime continues to resist a unified baseline once the
observational scaffold itself is hardened.

\section{Conclusion}

This paper is not written as a final declaration. We do not claim that ACM has
already conquered all low-acceleration systems, nor do we reduce MOND's
historical role to an erroneous episode. What the paper does establish is that
the interpretive authority of the supposed MOND stronghold can no longer rest
on a single label. It must instead be redistributed across distinct residual
mechanisms: geometry-locking, stellar-normalization bias, and low-structure
outer-gas regimes.

The main conclusion is therefore not ``who has already won.'' It is that
future discussions of the MOND-favored regime cannot be based on fit quality
alone. They must first pass through audits of distance, mass normalization,
outer structure, and source regime. The contribution of Paper II is to define
that objective audit framework.

This also means that harder distance anchors, more stable stellar mass
normalization, and higher-resolution outer-gas structure observations will do
more than supplement the present paper. They will continue to arbitrate the
interpretive ownership of these regimes. The hard31 in particular should not be
treated as a safe harbor for MOND, but as priority targets for future
higher-resolution observations and harder geometric distance determinations. We
do not pronounce a verdict in advance on behalf of future data. The present
audit nevertheless suggests that, as those data become harder, discussion of
the MOND-favored regime will become progressively less sustainable at the level
of surface fit quality alone and progressively more dependent on unified
baselines, residual pathology, and regime-level audit.

\section{Objective Validation Paths for Future Observations}

If the conclusions of this paper are to enter a genuinely reproducible phase,
future work should not remain at the level of swapping empirical fit curves. It
should instead proceed along three objectively higher-resolution paths. The
first is harder distance anchoring. Systems that remain in soft distance-source
regimes and exhibit strong step-like verdict responses under $D+e_D$ should be
prioritized for TRGB, Cepheid, maser, or comparable hard geometric distance
measurements~\cite{MadoreMagerFreedman2009}. The purpose is not to privilege
any theory in advance, but to stabilize the geometric coordinate system itself.

The second is more stable stellar mass normalization. Systems with low
$L_{3.6}$, high gas-to-light ratios, and strong sensitivity under conservative
$M/L$ scans should be re-examined under tighter stellar-population, color, and
evolutionary constraints. The purpose is to reduce the phase space that can be
absorbed by flexible normalization. The third is higher-resolution outer-disk
structure measurement. For hard31-like systems---low slope, low curvature, high
outer persistence, and no hard geometric anchor---the priority is not another
local corrective function. The priority is improved H\,I outer-disk
structure, more reliable inclination fields, and better constraints on
outer-disk asymmetry.

The paper therefore closes not with a demand that future observers accept a
particular theoretical conclusion, but with a validation program that any
position should accept. First harden the distance scale. Then stabilize the
stellar normalization. Then complete the outer-disk structural information.
Any theory that wishes to continue treating the MOND-favored regime as a
unified physical stronghold should first be required to pass through those same
higher-resolution audit channels.

\section{Data and Code Availability}

The theoretical baseline used in this paper is inherited from Paper~I, whose
preprint is available via \href{https://doi.org/10.5281/zenodo.19163963}{Zenodo DOI}
and whose public code repository is available at
\href{https://github.com/lesleyliuhaitao-ctrl/anchor-collapse-model.git}{anchor-collapse-model}. The
public release associated with the present paper, including the bilingual audit
frontend, the retained paper-II reproducible core, the curated evidence
package, and the archived audit-process scripts, is prepared at
\href{https://github.com/lesleyliuhaitao-ctrl/galaxy-audit-system.git}{galaxy-audit-system}. This
second repository serves as the public codebase for the audit layer developed
in Paper~II. Its interactive audit interface is publicly available at
\href{https://lesleyliuhaitao-ctrl.github.io/galaxy-audit-system/}{lesleyliuhaitao-ctrl.github.io/galaxy-audit-system/}
as a visual companion to the present manuscript.

\appendix

\section{Archived Audit Operators}

To prevent the main text from being overwhelmed by exploratory corrections,
we do not merge the full operator archive into the narrative body of the
paper. This omission should not be read as methodological irrelevance. On the
contrary, the hard31 subset acquires its status as a methodological judge
because a sequence of physically motivated, no-new-parameter operators was
tested against it and archived.

At minimum, this archive includes first-gradient corrections,
second-curvature corrections, diffusion-accumulation terms,
self-shielding terms, and nonlinear diffusion--shielding impedance
couplings. All of them obey the same retained-trunk discipline: no
per-galaxy freedom, no ad hoc rescue parameter, and no abandonment of the
unified baseline. In each case the question is simple. Can a physically
natural modification materially alter the hard31 verdict without compromising
the stability of the ACM-better majority?

The archive verdict is clearer than any one operator formula considered in
isolation. The tested operators were either too weakly activated to change
the hard31 classification, or they produced only marginal perturbations while
preserving global trunk safety. Their significance therefore lies not in
having nearly succeeded, but in forming a reproducible body of negative
results. This negative body of evidence moves the audit away from the repeated
search for one more elegant local function and toward the diagnosis of
historical bias on the input and regime sides.

In the final public form of the paper, the closed expressions, quantitative
summaries, and archive verdicts of these operators can be documented in a
longer appendix or linked research archive. For the present manuscript, it is
sufficient to note that the hard31 subset remained stable against a broad
family of no-new-parameter ACM-side rescue attempts. That stability is part
of the reason the subset functions as a methodological judge rather than a
generic collection of residual outliers.

\section{Supplementary Regime Diagnostics}

Not every diagnostic figure belongs in the main text. Some are more useful as
exclusionary evidence or regime-audit evidence than as headline figures. In
particular, the raw $V_{\mathrm{gas}}$ spectrum comparison between hard31 and
the ACM-better majority is retained primarily to clarify what the hard31
subset is \emph{not}. It is not merely a visibly fabricated frequency regime.
Nor is it reducible to a trivial smoothing artifact.

Likewise, the distance-family summary tables, ranked distance-edge response
lists, stellar-normalization sensitivity tables, and the hard31 geometry
audit summary remain important to reproducibility even when they do not serve
as the central visual anchors of the paper. Their function is not to amplify
the rhetoric of the main text. Their function is to keep the paper's
stronger claims traceable to explicit procedural audits. They can then be
checked item by item by external readers.

\bibliographystyle{apsrev4-2}
\bibliography{references}

\end{document}
